\documentclass[11pt]{article}
\usepackage[a4paper,margin=1in]{geometry}
\usepackage{amsmath,amssymb,amsthm,mathtools}
\usepackage[T1]{fontenc}
\usepackage{lmodern}
\usepackage{microtype}
\usepackage{enumitem}
\usepackage{booktabs}
\usepackage{hyperref}

\title{A Proposed LP Relaxation for the A Priori Traveling Salesman Path Problem}
\author{}
\date{}

\begin{document}
\maketitle

\section*{Problem setting}
Let $V$ be the set of customer vertices, let $r$ be a depot, and let $d$ be a metric on $\{r\} \cup V$. We consider an \emph{a priori traveling salesman path} problem: we first choose an ordering of the vertices in $V$, and then each vertex is activated independently with the same probability $p \in [0,1]$. Starting from the depot $r$, we visit the activated vertices in the chosen order and stop at the last activated vertex; there is no return to the depot.

For a permutation $\sigma$ of $V$, the random realized path starts at $r$ and visits the active vertices in the order
\[
\sigma(1),\sigma(2),\dots,\sigma(n), \qquad n := |V|,
\]
restricted to the activated set. The optimization goal is to minimize the expected cost of this realized path.

\section*{Decision variables}
We use the following variables.

\begin{itemize}[leftmargin=1.5em]
    \item $x_{v,i} \in [0,1]$ for each $v \in V$ and $i \in [n] := \{1,2,\dots,n\}$. Intuitively, $x_{v,i}$ is the weight that vertex $v$ occupies position $i$.
    \item $y_{v,i,u,j} \in [0,1]$ for each distinct positions $i,j \in [n]$ and distinct vertices $u,v \in V$. Intuitively, $y_{v,i,u,j}$ is the joint weight that $v$ is in position $i$ and $u$ is in position $j$.
    \item $t_i \ge 0$ for each $i=0,1,\dots,n$. Intuitively, $t_i$ is the expected arrival time at the vertex in position $i$, with the convention that if the vertex in position $i$ is inactive, then $t_i=t_{i-1}$.
\end{itemize}

\section*{Proposed LP relaxation}
The proposed LP is:

\begin{alignat}{2}
\min \quad & t_n \label{obj}\\
\text{s.t.}\quad
& \sum_{i=1}^n x_{v,i} = 1 && \forall v \in V, \label{assign-vertex}\\
& \sum_{v \in V} x_{v,i} = 1 && \forall i \in [n], \label{assign-position}\\
& \sum_{u \in V \setminus \{v\}} y_{v,i,u,j} = x_{v,i} && \forall v \in V,\ \forall i,j \in [n],\ i \neq j, \label{pair-left}\\
& \sum_{v \in V \setminus \{u\}} y_{v,i,u,j} = x_{u,j} && \forall u \in V,\ \forall i,j \in [n],\ i \neq j, \label{pair-right}\\
& \sum_{v \in V} \sum_{u \in V \setminus \{v\}} y_{v,i,u,j} = 1 && \forall i,j \in [n],\ i \neq j, \label{pair-total}\\
& t_0 = 0, \label{t0}\\
& t_i = t_{i-1} + p(1-p)^{i-1}\sum_{v \in V} d(r,v)x_{v,i} 
   && \forall i \in [n], \label{ti-main}\\
& \hspace{4em} + \sum_{j=1}^{i-1} p^2(1-p)^{i-j-1}
   \sum_{v \in V}\sum_{u \in V \setminus \{v\}} d(u,v)\, y_{v,i,u,j}
   && \forall i \in [n], \notag\\
& x_{v,i} \ge 0 && \forall v \in V,\ \forall i \in [n], \label{xnonneg}\\
& y_{v,i,u,j} \ge 0 && \forall u,v \in V,\ u \neq v,\ \forall i,j \in [n],\ i \neq j. \label{ynonneg}
\end{alignat}

\paragraph{Remark on redundancy.}
Constraint \eqref{pair-total} is implied by \eqref{pair-left} and \eqref{assign-position}, or equivalently by \eqref{pair-right} and \eqref{assign-position}. It is included explicitly because it is a natural local-consistency condition and sometimes makes the formulation easier to read.

\section*{Interpretation of the recurrence}
For a fixed position $i$, the term
\[
p(1-p)^{i-1}\sum_{v \in V} d(r,v)x_{v,i}
\]
represents the expected contribution when the vertex in position $i$ is the first active vertex among positions $1,2,\dots,i$.

For each $j < i$, the term
\[
p^2(1-p)^{i-j-1}
\sum_{v \in V}\sum_{u \in V \setminus \{v\}} d(u,v)\, y_{v,i,u,j}
\]
represents the expected contribution when the vertex in position $j$ is the most recent active predecessor of the vertex in position $i$, that is:
\begin{itemize}[leftmargin=1.5em]
    \item the vertex at position $j$ is active,
    \item the vertex at position $i$ is active, and
    \item all positions strictly between $j$ and $i$ are inactive.
\end{itemize}
The probability of this event is exactly $p^2(1-p)^{i-j-1}$ under identical independent activation.

\section*{Deterministic special case: $p=1$}
When $p=1$, the recurrence collapses to
\[
t_1 = \sum_{v \in V} d(r,v)x_{v,1},
\]
and for every $i \ge 2$,
\[
t_i = t_{i-1} + \sum_{v \in V}\sum_{u \in V \setminus \{v\}} d(u,v)\, y_{v,i,u,i-1}.
\]
Thus, for an integral solution corresponding to a permutation $\sigma$,
\[
t_n = d(r,\sigma(1)) + \sum_{i=2}^n d\bigl(\sigma(i-1),\sigma(i)\bigr),
\]
which is exactly the cost of the Hamiltonian path starting at the depot and following the ordering $\sigma$ without returning to $r$.

\section*{Discussion}
This formulation should be viewed as a \emph{proposed LP relaxation}. The variables $x$ encode first-order position marginals, and the variables $y$ encode pairwise position marginals together with local consistency constraints. Whether these local constraints are sufficient for a desired global lower-bound property (for example, domination of the minimum spanning tree cost in the deterministic case) requires separate analysis.

\end{document}
